\section{Matching elementary sides under the actual subdivision}
\label{sec:area_law_actual_side_matching}

When two actual fine cells have a common side segment, their midpoint
choices must agree on that segment. The following statement applies to
each elementary segment under the subdivision determined by all actual
corners at or above the fixed lower index; see
\cite[Section~11]{OpenAI2026AreaLaw}.

% Source: 10-geometry.tex, lines 299--306 and 352--356,
% prop:two-families and geometry:initial-stars,
% revision adc7f1241b42e322a6451854ab7e4b4c146bf78a.
% This is conditional matching, not opponent existence or a region assignment.

\begin{theorem}[Every elementary segment with contact matches]
    \label{thm:area_law_actual_elementary_side_matching}
    \lean{TNLean.PEPS.AreaLaw.Geometry.fineLayer_elementary_contact_match}
    \leanok
    \uses{def:area_law_actual_side_mask,def:area_law_fine_layer_indices,
        def:area_law_fine_belt_scales}
    Fix a common origin $o\in\mathbb R^2$, finite endpoint set $Z$,
    integer width $C_0\ge2$ and natural indices $k_0,k,h$ with
    $k\ge50\,000\,000$ and $k_0\le k,h$.
    Let $z\in F_{k,\ell_k}$ and $w\in F_{h,\ell_h}$ be actual
    fine-cell indices with $(k,z)\ne(h,w)$, and write
    $Q_{k,z}=Q_{\ell_k}(z)$ and $Q_{h,w}=Q_{\ell_h}(w)$.
    Let $[a_i,b_i]$ be any elementary segment of the reference cell
    under its actual midpoint subdivision relative to $k_0$.
    Suppose that
    $[a_i,b_i]\cap\overline{Q_{h,w}}$ contains two distinct points.
    Then there is an elementary segment $[a_j,b_j]$ of the other
    cell under its actual midpoint subdivision relative to $k_0$
    such that
    \begin{align}
        s(j) & =s(i)+2\pmod4,       \\
        a_i  & =b_j,\qquad b_i=a_j,
    \end{align}
    where $s(i)$ and $s(j)$ are their counterclockwise side indices.
\end{theorem}
\begin{proof}\leanok
    \uses{thm:area_law_fine_cell_contact,thm:area_law_elementary_side_geometry,
        def:area_law_actual_side_mask,thm:area_law_elementary_corner_endpoint,
        thm:area_law_unsplit_side_endpoints,thm:area_law_adjacent_fine_meshes}
    The closed-cell contact is either a whole side of each cell or
    a whole side of the smaller cell equal to a midpoint half of a
    side of the larger cell. Contact containing two distinct points
    identifies the facing side of the reference cell: two different
    sides of the same square meet in at most one point.

    In the whole-side case, the two whole sides have the same
    midpoint. Their subdivision conditions therefore have exactly
    the same corner witnesses. If they are undivided, the whole
    sides match. If they are divided, the first half on one side
    matches the second half on the opposite side, and conversely.

    In the unequal-size case, a corner of the smaller cell is the
    midpoint of the larger side. The smaller cell lies at or above
    $k_0$, so this corner forces the larger side to be divided.
    The smaller side is undivided. Indeed, a corner at its midpoint
    would lie strictly inside the corresponding elementary half
    of the larger side, contrary to the elementary-corner endpoint
    assertion. That assertion is applicable to the larger cell:
    its layer is at least that of the smaller cell, so the required
    late-layer bound follows from the reference bound whenever the
    reference cell is the smaller one.

    Hence the whole smaller side matches the appropriate half of
    the larger side. When the reference cell is larger, contact
    containing two distinct points selects that half, since the
    two halves meet only at their midpoint. Finally, counterclockwise
    traversal of opposite facing sides gives the reversed endpoint
    equalities.
\end{proof}

Theorem~\ref{thm:area_law_elementary_side_opponent} supplies an opposing
closed region for every elementary segment. Opponent uniqueness,
consistent region labels and the isolated-star construction remain
further steps in the two-family argument.
